\chapter{Trò chuyện qua không khí}
% full version: chapters/05_serocomputer.tex

\pencilfig{5}{\pencilart{5}}{Bên trái là điện thoại: dây dẫn từ một tới một. Bên phải là radio: một thông điệp gửi cùng lúc tới mọi thứ xung quanh.}

Cho đến giờ các nơ-ron nói chuyện với nhau qua điện thoại: dây nối dây, có địa chỉ. Giờ hãy bật đài phát thanh đầu tiên --- các nơ-ron \nt{SERO}, nhả serotonin. Cả bộ não chỉ có vài trăm nghìn nơ-ron như vậy, nhưng dây dẫn của chúng tỏa ra khắp những vùng não rộng lớn.

\section*{Bốn điểm khác biệt}

Nơ-ron \nt{SERO} không giống \nt{GLUT}. Thứ nhất, dấu của tín hiệu do nơi nhận chọn: serotonin có ``ổ khóa'' của cả hai dấu. Thứ hai, nó tự hoạt động: khi thức, nó tách đều đặn vài lần mỗi giây, nếu nhận được một nguồn nuôi nền ổn định --- nó là một chiếc máy đếm nhịp. Thứ ba, nó chậm: tín hiệu của nó tác động không trong vài mili giây mà trong hàng trăm mili giây. Thứ tư, nó thường nhả chất dẫn truyền không vào khe hẹp mà ra khoảng không xung quanh, nơi chất này loang ra như giọt mực trong nước. Nơi nhận cảm nhận được nồng độ và không biết phân tử đến từ ai.

\section*{Thứ logic không tồn tại}

Máy đếm nhịp có ``ổ khóa'' ức chế là một cổng ``không'' làm sẵn: khi chưa có serotonin, nơ-ron hoạt động; serotonin xuất hiện --- nó im lặng. Một nơ-ron thay cho cả một mạch. Có vẻ như mạng serotonin mạnh hơn mạng \nt{GLUT} về logic.

Nhưng có một vướng mắc. Chất dẫn truyền đã nhả ra khoảng không sẽ trộn lẫn. Hai tín hiệu chỉ còn khác nhau nếu chúng được nhả ở những điểm cách nhau xa hơn quãng đường chất kịp loang ra --- tức vài chục micromet. Mà mỗi nơ-ron serotonin lại có vô số điểm nhả, rải khắp những vùng rộng, và ``dây dẫn'' của các nơ-ron khác nhau chồng lên nhau. Không thể truyền từng bit riêng lẻ qua một mạng như vậy. Nó chỉ có thể truyền vài đại lượng chung --- như chiếc radio mà mọi người cùng nghe một đài.

\pencilfig{123}{%
% two release sites: far apart the clouds stay separate (two signals), close together they merge (one level)
\foreach \x/\y in {0.9/1.55, 4.1/1.55, 1.9/0, 2.9/0}{
  \foreach \r/\o in {0.95/0.07, 0.7/0.09, 0.45/0.12}{\fill[pcViolet, opacity=\o] (\x,\y) circle (\r);}
  \draw[dotpen=pcViolet, fill=pcViolet] (\x,\y) circle (0.07);}
\draw[arr] (1.95,1.55) -- (3.05,1.55); \draw[arr] (3.05,1.55) -- (1.95,1.55);
\node[lbl, anchor=south, align=center] at (2.5,1.6) {xa hơn\\độ loang};
\node[lbl, anchor=west] at (5.25,1.55) {hai tín hiệu};
\node[lbl, anchor=west] at (5.25,0) {một đám mây chung};
}{Serotonin được nhả ra khoảng không và loang thành một đám mây. Hai nguồn chỉ cho hai tín hiệu khác nhau nếu chúng cách nhau xa hơn quãng chất kịp loang ra. Nếu gần hơn, các đám mây nhập thành một mức chung.}

\section*{Nó thực sự làm gì}

Để giữ một đại lượng chung, hệ thống này có đủ mọi thứ. Nồng độ serotonin làm mượt từng tiếng tách riêng lẻ thành một mức trơn tru --- giống như bộ chuyển đổi đơn giản nhất biến dòng xung thành điện áp. Mức này giữ được vài giây rồi tự tan: không phải một bit cần xóa, mà là một dấu vết mờ dần.

Nơ-ron serotonin còn có ``ổ khóa'' trên chính mình: serotonin nhả ra hãm bớt chính nguồn của nó. Kỹ sư sẽ nhận ra ở đây một bộ khuếch đại có phản hồi âm --- bộ ổn nhiệt giữ mức. Có điều, thí nghiệm cho thấy sự tự ức chế này có địa chỉ và chỉ gây ra những khoảng ngừng ngắn, nên vai trò bộ ổn nhiệt hiện vẫn là giả thuyết hơn là sự thật.

\section*{Chờ đợi đáng giá bao nhiêu}

Đại lượng chung nào nên giữ chậm và như nhau cho toàn mạng? Một ứng viên tốt là \emph{sự kiên nhẫn}. Hãy hình dung một lựa chọn: phần thưởng nhỏ ngay bây giờ hay phần thưởng lớn nhưng sau vài giây. Bạn sẵn lòng chờ bao lâu phụ thuộc vào việc tương lai ``mất giá'' nhanh thế nào đối với bạn. Theo một giả thuyết, chính serotonin vặn núm này. Có những thí nghiệm ủng hộ nó: nếu bật nhân tạo các nơ-ron serotonin, con chuột chờ phần thưởng bị hoãn lâu hơn trước khi bỏ cuộc.

Chương này giới thiệu ba thiết bị mà mọi đài phát thanh của bộ não sẽ dùng: nơ-ron máy đếm nhịp, ``dây dẫn'' thực chất là một thể tích, và một mức chậm thay cho từng tiếng tách riêng lẻ.

\fullversion{5}
